% ======================================================================
%  PARTIE 1 -- Contexte et objectifs                        (~2 min 15)
% ======================================================================
\section{Contexte et objectifs}

% ----------------------------------------------------------------------
% À DIRE (~45 s)
%  - Stage de 6 mois au laboratoire CRIStAL (Centre de Recherche en
%    Informatique, Signal et Automatique de Lille), UMR 9189 sous tutelle
%    du CNRS, de l'Université de Lille, de Centrale Lille et d'Inria.
%  - Équipe ToSyMA (Tolérance aux fautes des Systèmes Mobiles Autonomes) :
%    diagnostic, observation d'état, fusion multi-capteurs tolérante aux
%    fautes pour la localisation -> lien direct avec mon sujet.
%  - Projet ANR-SOS : doter des flottes de robots mobiles, en extérieur,
%    de capacités de cartographie et de coopération, même quand la
%    communication entre les robots n'est pas continue.
%  - Ma mission : la brique "cartographie sémantique" : construire une
%    carte enrichie sémantiquement à partir du LiDAR et de la caméra d'un
%    robot, puis étendre cette brique au multi-robots.
% ----------------------------------------------------------------------
\begin{frame}{Contexte du stage}
  \vfill
  \begin{tcbraster}[raster columns=3, raster equal height, raster column skip=0.45cm]
    \begin{tcolorbox}[carte, title=Laboratoire]
      \alert{CRIStAL} -- Lille\\
      {\footnotesize UMR 9189}\\[8pt]
      Équipe \alert{ToSyMA}\\
      {\footnotesize\color{textgray} fusion multi-capteurs\\ tolérance aux fautes}
    \end{tcolorbox}
    \begin{tcolorbox}[carte, title=Projet ANR-SOS]
      Flottes de \alert{robots mobiles}\\
      en \alert{extérieur}\\[8pt]
      Cartographie \& coopération\\
      {\footnotesize\color{textgray} communication imparfaite\\ entre les robots}
    \end{tcolorbox}
    \begin{tcolorbox}[carte, title=Mon stage]
      Brique de\\ \alert{cartographie sémantique}\\[8pt]
      Fusion \alert{LiDAR -- caméra}\\[4pt]
      Extension \alert{multi-robots}
    \end{tcolorbox}
  \end{tcbraster}
  \vfill
  \centering
  \legende{Stage du 27 mars au 27 septembre 2026 \quad\textbullet\quad Encadrante : Cindy Cappelle (CRIStAL)}
\end{frame}

% ----------------------------------------------------------------------
% À DIRE (~45 s)
%  - Une carte purement géométrique (nuage de points, occupation) dit OÙ
%    sont les obstacles, mais pas CE QUE c'est.
%  - La carte sémantique ajoute une étiquette à chaque élément (route,
%    bâtiment, végétation, véhicule, piéton...) : information directement
%    exploitable pour la navigation autonome et la décision de haut
%    niveau (ex. privilégier une route plutôt qu'une zone de végétation).
%  - Dans ANR-SOS, cette carte doit en plus être partagée et fusionnée
%    entre plusieurs robots dont la communication est imparfaite.
%  - D'où la question : comment construire une carte sémantique partagée
%    entre robots à partir du LiDAR (géométrie précise) et de la caméra
%    (information sémantique) ?
%  - (slam.jpg est une image d'illustration : citer la source si elle ne
%    vient pas de tes propres acquisitions.)
% ----------------------------------------------------------------------
\begin{frame}{Problématique}
  \centering
  \begin{columns}[T]
    \begin{column}{0.48\textwidth}
      \centering
      \includegraphics[height=3.25cm]{slam.jpg}\\[4pt]
      \textbf{Carte géométrique} $\rightarrow$ \alert{où} ?\\
      \legende{nuage de points, obstacles}
    \end{column}
    \begin{column}{0.48\textwidth}
      \centering
      \includegraphics[height=3.25cm]{points.png}\\[4pt]
      \textbf{+ Sémantique} $\rightarrow$ \alert{quoi} ?\\
      \legende{route, bâtiment, végétation, véhicule\dots}
    \end{column}
  \end{columns}
  \vspace{0.3cm}
  \begin{tcolorbox}[question, width=\textwidth]
    Comment construire une carte \alert{sémantique}, \alert{partagée} entre robots,\\
    à partir du \alert{LiDAR} et de la \alert{caméra} ?
  \end{tcolorbox}
\end{frame}

% ----------------------------------------------------------------------
% À DIRE (~45 s)
%  - Démarche : deux briques existantes, complémentaires :
%      * LIO-SAM (LiDAR + IMU) : localisation précise du robot et carte
%        géométrique -> mais aucune information sémantique ;
%      * YOLO (caméra) : segmentation sémantique de l'image -> mais pas de
%        géométrie 3D.
%  - Mon travail = le framework ROS 2 qui les fusionne : chaque point
%    LiDAR reçoit un label YOLO, puis est placé dans le repère monde grâce
%    à la pose LIO -> carte sémantique 2D/3D exportée en GeoJSON.
%  - Puis trois extensions : prise en compte des loop closures, carte a
%    priori (OpenStreetMap), et passage au multi-robots.
%  - Validation : jeu de données KITTI, puis deux robots AgileX réels
%    (Scout Mini et Ranger).
% ----------------------------------------------------------------------
\begin{frame}{Objectifs et démarche}
  \centering
  \begin{tikzpicture}[font=\scriptsize]
    % Capteurs
    \node[capteur] (lidar) at (0,1.0)   {LiDAR 3D};
    \node[capteur] (imu)   at (0,0.05)  {IMU {\tiny(+ GPS)}};
    \node[capteur] (cam)   at (0,-1.25) {Caméra};
    % Briques existantes
    \node[boite=2.6cm, minimum height=0.95cm] (lio)  at (2.85,0.52)  {\textbf{LIO-SAM}\\ pose + carte géométrique};
    \node[boite=2.6cm, minimum height=0.95cm] (yolo) at (2.85,-1.25) {\textbf{YOLO}\\ segmentation sémantique};
    % Contribution
    \node[donnee=2.4cm, minimum height=1.2cm] (fus) at (6.0,-0.35) {\textbf{Fusion}\\ LiDAR -- caméra\\ {\tiny points 3D étiquetés}};
    \node[sortie=2.3cm, minimum height=1.2cm]  (carte) at (8.95,-0.35) {Carte sémantique\\ 2D / 3D\\ {\tiny\mdseries GeoJSON}};
    \node[boite=2.3cm, fill=softblue, minimum height=1.2cm] (multi) at (11.85,-0.35) {\textbf{Multi-robots}\\ carte commune};
    % Flèches
    \draw[fleche] (lidar.east) -- ++(0.3,0) |- (lio.west);
    \draw[fleche] (imu.east) -- ++(0.3,0) |- (lio.west);
    \draw[fleche] (cam) -- (yolo);
    \draw[fleche] (lidar.north) -- ++(0,0.35) -| node[pos=0.25, above, text=textgray, font=\tiny] {nuage de points} (fus.north);
    \draw[fleche] (lio.east) -- ++(0.35,0) |- (fus.west);
    \draw[fleche] (yolo.east) -- ++(0.35,0) |- (fus.west);
    \draw[fleche] (fus) -- (carte);
    \draw[fleche] (carte) -- (multi);
    % Ce qui existe / ce qui est développé
    \begin{scope}[on background layer]
      \node[groupe, fit=(lio)(yolo)] (existant) {};
      \node[groupe, draw=accent!70, fit=(fus)(carte)(multi)] (stage) {};
    \end{scope}
    \node[font=\tiny, text=textgray, anchor=north] at (existant.south) {briques existantes};
    \node[font=\tiny\bfseries, text=accent, anchor=north] at (stage.south) {développé pendant le stage};
  \end{tikzpicture}

  \vspace{0.45cm}
  {\small + \alert{loop closure} \quad\textbullet\quad + \alert{carte a priori} (OpenStreetMap)\\[3pt]
   Validation : \alert{KITTI} $\rightarrow$ robots réels \alert{Scout Mini} et \alert{Ranger}}
\end{frame}
